\documentclass[12pt,letterpaper]{article}
\usepackage{amsmath, amssymb}
\usepackage[dvips]{graphicx}
\usepackage{graphicx}
\usepackage{epsfig, float}
\usepackage[rflt]{floatflt}
\usepackage{multicol}
\usepackage{ textcomp,wasysym }
\usepackage{pstricks-add}
\usepackage{rotating}
\usepackage[T1]{fontenc}
\usepackage{calligra}
\usepackage{pbsi}
\usepackage{longtable}
\topmargin -0.5 in \textwidth 6.8in \textheight 9in \oddsidemargin
-0.5cm

\begin{document}
\title{Lab Numerico: calculo computacional}
{\begin{minipage}{2cm}
\hspace*{1cm}\includegraphics[width=0.6\textwidth]{escubo-ubb.eps}
\end{minipage}
\begin{minipage}{12cm}
\small
{\bf \rm
{\hspace*{3.1cm} {\footnotesize UNIVERSIDAD DEL B\'IO-B\'IO} \\
 \hspace*{3.7cm} {\scriptsize FACULTAD DE CIENCIAS}  \\
  \hspace*{2.9cm} {\scriptsize DEPARTAMENTO DE MATEM\'ATICA}  \\
\hspace*{3.75cm} {\scriptsize Autor: Franco Milanese}\\
\hspace*{3.75cm} {\scriptsize Segundo Semestre de 2014}}}
\end{minipage}}
{\begin{minipage}{2cm}
\hspace*{-0.5cm}\vspace*{-0.05cm}\includegraphics[width=0.7\textwidth]{escudo-dmat.eps}
\end{minipage}}

\hspace*{-1,5cm}\rotatebox[origin=c]{90}{\begin{picture}(0,0)
\put(0,7){\makebox(9,-13)[l]{\hspace*{-6.5in} \bf \it Departamento de Matem\'atica - Universidad del B\'io-B\'io - 2014}}
\end{picture}}

\vspace*{0.1cm} \centerline {\bf\underline{Laboratorio N 3 C\'alculo Num\'erico (220098 \& 220023): C\'alculo Computacional}} \vskip0.5cm


\section{Sobre los errores}

Resolver un problema num\'erico es casi sin\'onimo de obtener una cierta aproximaci\'on a
un cierto valor num\'erico $x_e$. Si $x_a$ es una de esas aproximaciones al valor $x_e$ , se define el
error absoluto como
$$
abs(x_e,x_a)=|x_e - x_a|.
$$
Para muchos prop\'ositos, sin embargo, es preferible considerar el porcentaje de error o
error relativo
$$
rel(x_e,x_a)=\frac{|x_e-x_a|}{|x_e|} \quad \text{ si } |x_e|\neq 0
$$
En t\'erminos generales, la utilidad de manejar el error relativo radica en que nos protege
contra afirmaciones precipitadas sobre la bondad de una aproximaci\'on, sobre todo cuando
nos movemos en escalas extremas (n\'umeros muy grandes o muy peque\~{n}os).

El concepto de error relativo est\'a relacionado con la representaci\'on decimal a trav\'es de
la noci\'on de d\'igitos significativos correctos. En concreto, puede afirmarse que si
$$
rel(x_e,x_a) \leq  0.5 \times 10^{-m}, \quad m \in \mathbb{N}
$$
entonces, $x_a$ comparte $m$ d\'igitos significativos correctos con $x_e$ despu\'es de redondear ambos
valores a $m$ cifras significativas

\section{ \textbf{Sobre la precisi\'on y los formatos de n\'umeros}}

Casi todas las operaciones en \texttt{MATLAB} son realizadas con precisi\'on \texttt{double} de acuerdo al est\'andar 754 de la IEEE. Debido a que los computadores representan n\'umero con una precisi\'on finita, los c\'alculos pueden llevar a resultados que matem\'aticamente no son intuitivos. Es importante resaltar que estos resultados no son errores de MATLAB sino mal uso de los usuarios.

\texttt{MATLAB} incluye varios formatos de precisi\'on, entre ellos 

\begin{longtable}{|c||p{0.7\textwidth}|}
\hline
double	& double(x) es una funci\'on sobrecargada que convierte a precisi\'on doble todo objeto que tenga sentido convertir. \\
single 	& single(A) convierte la matriz A a precisi\'on single. A puede ser cualquier objeto num\'erico como un double. Los n\'umero con precisi\'on single requieren menos espacio que los de precisi\'on doble. Pero tienen menos precisi\'on y un rango menor.\\
int8 	& int8(array) convierte los elementos de un arreglo en n\'umero de 8-bit (1-byte)  con signo $\pm 2^7-1$. \\	
& an\'alogamente 
int16	
int32
int64\\
uint8 & uint8(array) convierte los elementos de un arreglo en n\'umero de 8-bit (1-byte)  sin signo $\pm 2^8-1$. \\	
& an\'alogamente 
uint16
uint32
uint64\\
\hline
\end{longtable}

El siguiente c\'odigo ejemplifica distintos tipos de formatos en una implementaci\'on en MATLAB

\begin{verbatim}
a = magic(4);
b = single(a);

whos
  Name      Size         Bytes  Class

  a         4x4            128  double array
  b         4x4             64  single array
\end{verbatim}

MATLAB tambi\'en permite verificar el formato de sus variables con las siguientes funciones 

\begin{longtable}{|c||p{0.7\textwidth}|}
\hline
isinteger 	&	isinteger(A) retorna TRUE si el arreglo A es de alg\'un tipo entero, retorna FALSO  si no. \\
isfloat 	& 	isfloat(A) retorna TRUE si el arreglo A es de formato punto flotante. Los tipos de formatos punto flotante son single() y double(). \\
isnumeric	& 	isnumeric(A) retorna verdad si A es del tipo single, double, int8, uint8, int16, uint16, int32, uint32, int64, uint64 \\
isreal 		& 	isreal(A) retorna TRUE si A no tiene partes imaginarias. \\
isfinite	& 	isfinite(A)  retorna un arreglo del mismo tama\~{n}o que A que contiene unos l\'ogicos donde A es finito y ceros l\'ogicos donde son infinitos o NaN. \\
isinf 		& 	isinf(A)  retorna un arreglo del mismo tama\~{n}o que A que contiene unos l\'ogicos donde A es +Inf o -Inf y ceros l\'ogicos donde no. \\
isnan		& 	isnan(A)  retorna un arreglo del mismo tama\~{n}o que A que contiene unos l\'ogicos donde A es Nan y ceros l\'ogicos donde no.\\
class()		& 	class(A) entrega el tipo de variable que es A en una cadena. \\
isa()		&   is(A,'tipo') entrega el valor de verdad de la proposici\'on class(a)=='tipo'
		\\
\hline
\end{longtable}

Para cualquier real A, s\'olo una de las cantidades isfinite(A), isinf(A) y isnan(A) es igual a 1.

\begin{verbatim}
isinteger(uint8(1:255))
ans =
     1

isinteger(pi)
ans =
     0

isinteger(3)
ans =
     0

isinteger(isinteger(uint8(3)))
ans =
     0

A = [-2  -1  0  1  2];

isnan(1./A)
ans =
     0     0     0     0     0

isnan(0./A)
ans =
     0     0     1     0     0
\end{verbatim}

En cada formato MATLAB es capaz de determinar su precisi\'on relativa usando la funci\'on eps().

\begin{longtable}{|c||p{0.7\textwidth}|}
\hline
eps		& 
Usos


eps

d = eps(X)

eps('double')

eps('single')

Descripci\'on

eps retorna la distancia desde 1.0 hasta el pr\'oximo n\'umero de precisi\'on double. (eps=$2^{52}$)

d = eps(X) es la distancia desde abs(X) hasta el pr\'oximo entero como punto flotante que tenga la misma precisi\'on que X. X puede ser de precisi\'on double o single.
\\
\hline
\end{longtable}

Funciones valores cr\'iticos en el an\'alisis de error de un programa seg\'un precisi\'on.

\begin{longtable}{|c||p{0.7\textwidth}|}
\hline
flintmax		& 	flintmax() retorna el mayor entero consecutivo que puede almacenar con precisi\'on double ($2^{53}$). Sobre este valor, la precisi\'on double() no tiene precisi\'on entera y no todos los enteros se pueden representar exactamente. \\
intmax			& 	intmax() es el mayor valor positivo que se puede representar en MATLAB con un entero de 32 bits. Cualquier valor mayor que intmax() se trunca a intmax() cuando es convertivo a int32. \\
intmin			& 	intmin() es el menor valor se puede representar en MATLAB con un entero de 32 bits. Cualquier valor menor que intmin() se trunca a intmin() cuando es convertivo a int32. \\
realmax			& 	realmax() retorna el mayor n\'umero en formato double. \\
realmin			& 	realmin() retorna el menor n\'umero en formato double. \\
\hline
\end{longtable}

\begin{verbatim}
str = 'El rango de los double es :\n\t%g hasta  %g y \ n\t %g hasta  %g';
sprintf(str, -realmax, -realmin, realmin, realmax)

ans =
El rango de los double es
   -1.79769e+308 hasta  -2.22507e-308 y
    2.22507e-308 hasta  1.79769e+308
\end{verbatim}


%%%%%%%%%%%%%%%%
%%%%%%%%%%%%%
\section{\textbf{Ejemplos de errores de precisi\'on y formato}}

\begin{enumerate}

\item 
\begin{verbatim}
format long

eps(5)
ans =
     8.881784197001252e-16
\end{verbatim}

Estas instrucciones muestran que no hay n\'umeros de precisi\'on double() entre 5  y 5+eps(5). Si se ejecuta un c\'alculo con precisi\'on double 
que retorna 5, este resultado tiene una precisi\'on de eps(5).

El valor de eps(x) depende de x. A medida que x aumenta, tambi\'en lo hace su precisi\'on:

\begin{verbatim}
eps(50)
ans =
     7.105427357601002e-15
\end{verbatim}

\item  Para un n\'umero x del tipo double(), eps(single(x)) entrega una cota superior para el valor al que x es redondeado cuando es convertido de double a single. Por ejemplo

\begin{verbatim}
double(single(3.14) - 3.14)
ans =
   1.0490e-07
\end{verbatim}
El valor al que $3.14$ es redondeado es menos que 
\begin{verbatim}
eps(single(3.14))
ans =
  2.3842e-07
\end{verbatim}


\item El n\'umero irracional $4/3$ no es representable exactamente como una divis\'on en n\'umeros binarios. Por esta raz\'on el siguiente c\'alculo no es cero, pero muestra la cantidad eps()

\begin{verbatim}
>>e = 1 - 3*(4/3 - 1)

e =
   2.2204e-16
\end{verbatim}

Similarmente, \texttt{0.1} no es representable exactamente como binario, en consecuencia, sucede que 

\begin{verbatim}
>>a = 0.0;
for i = 1:10
  a = a + 0.1;
end

>>a == 1

ans =
     0
     
\end{verbatim}

\item Hay saltos entre n\'umeros en punto flotante, a medida que los n\'umeros crecen estos saltos tambi\'en, esto queda evidencia en el siguiente c\'odigo

\begin{verbatim}
(2^53 + 1) - 2^53

ans =
     0
\end{verbatim}

Como \texttt{pi} no es realmente $\pi$, no es sorprendente que 

\begin{verbatim}
sin(pi)

ans =
     1.224646799147353e-16
\end{verbatim}

\item Cuando se realizan restas con t\'erminos muy parecidos, pueden suceder cancelaciones inesperadas. El siguiente ejemplo 
muestra una cancelaci\'on por p\'erdida de la precisi\'on, lo que hace la suma insignificante. 

\begin{verbatim}
sqrt(1e-16 + 1) - 1

ans =
     0
\end{verbatim}

\item El siguiente ejemplo muestra como dibujar un polinomio de grado 7.

\begin{verbatim}
x = 0.988:.0001:1.012;
y = x.^7-7*x.^6+21*x.^5-35*x.^4+35*x.^3-21*x.^2+7*x-1;
plot(x,y)
\end{verbatim}

El resultado se se obtiene no aparenta ser un polinomio. La gr\'afica no es para nada suave. Lo 
que se muestran son errores de redondeos en acci\'on. La escala del eje $y$ es pequeña, estos pequeños valores de $y$ 
son calculados usando n\'umeros tan grandes como $35 \cdot 1,012^4$  y adem\'as hay varias cancelaciones en sustracciones.

\end{enumerate}

\section{\textbf{Ejercicios}}

\begin{enumerate}
\item Ingrese las instrucciones 
\begin{verbatim}
A = [1 1; 1 0]
X = [1 0; 0 1]
\end{verbatim}
luego 
\begin{verbatim}
X = A*X
\end{verbatim}
Ahora, presione repetidas veces la flecha hacia arriba seguida de la tecla Enter. \textquestiondown Qu\'e sucede ?, \textquestiondown Reconoce 
las matrices que est\'an siendo generadas?, \textquestiondown Cuantas veces podra hacer esto antes de que suceda un overflow?.

\item \textquestiondown Qu\'e hacen los siguientes c\'odigos?, \textquestiondown Cuantas l\'ineas de salidas producen?, \textquestiondown Cuales son los dos \'ultimos valores de $x$ que se imprimen?.

\begin{verbatim}
x = 1; while 1+x > 1, x = x/2, pause(.02), end

x = 1; while x+x > x, x = 2*x, pause(.02), end

x = 1; while x+x > x, x = x/2, pause(.02), end
\end{verbatim}

\item El desarrollo en serie de Taylor del seno es 
$$sen(x) = x-\frac{x^3}{3!}+\frac{x^5 }{5!}- \cdots$$
La siguiente funci\'on de Matlab aproxima la funci\'on seno 
usando su desarrollo en serie de Taylor.
\begin{verbatim}
function s=powersin(x)
% POWERSIN. Power series for sin(x).
% POWERSIN(x) tries to compute sin(x)
% from a power series
s = 0;
t = x;
n = 1;
while s+t ~= s;
s = s + t;
t = -x.^2/((n+1)*(n+2)).*t;
n = n + 2;
end
\end{verbatim}
\begin{enumerate}
\item \textquestiondown Porqu\'e se detiene el ciclo while?.
\item Cuando $x\in\{\frac{\pi}{2}, 11\frac{\pi}{2}, 21\frac{\pi}{2}, 31\frac{\pi}{2}\}$, 
\textquestiondown Cuan preciso es el resultado calculado por \texttt{powersin(x)}?, \textquestiondown Cu\'antos t\'erminos del desarrollo en potencia son utilizados?, \textquestiondown Cu\'al es el t\'ermino mas grande de la serie?.
\item \textquestiondown Que concluye de la aritm\'etica de punto flotante y las series de potencias para evaluar funciones?.
\end{enumerate}

\item Explique la salida producida por
\begin{verbatim}
>>t = 0.1
n = 1:10
e = n/10 - n*t
\end{verbatim}


\item 

\begin{verbatim}
>> v = sin(2*pi)==sin(4*pi)

v =

     0
\end{verbatim}

\begin{enumerate}
\item Determine cual es el origen de este error.

\item Proponga otra forma de verificación que permita asegurar el comportamiento que anal\'iticamente se espera de la instrucci\'on anterior.

\end{enumerate}

\item Un famoso resultado de L. Euler (1707-1783) afirma que
$ \displaystyle \sum_{n=1}^\infty \frac{1}{n^2} = \frac{\pi^2}{6}$.
Use el n\'umero de precisi\'on double 
\begin{verbatim}
>> pi^2/6
\end{verbatim}
como una representaci\'on exacta, para obtener 
el menor número de sumandos de la serie anterior, de modo que la correspondiente
suma finita aproxime a $\frac{\pi^2}{6}$ con un error menor o igual que \texttt{10-6}.

\end{enumerate}


\begin{thebibliography}{9}

\bibitem{Mo1} Moler, Cleve, Floating Points, MATLAB News and Notes, Invierno, 1996. 

\bibitem{Mo2} Moler, Cleve, Numerical Computing with MATLAB, S.I.A.M. Disponible en http://www.mathworks.com/moler/.

\bibitem{H1} N. J. Higham, Accuracy and Stability of Numerical Algorithms, SIAM, Philadelphia, 2002
\end{thebibliography}

\end{document}